\chapter*{Abstract}

The evolution of deep learning has enabled human-dependent work to be handled by models, especially in image segmentation. The segmentation problem in computer vision has been applied in various fields and has achieved successes due to the many architectures and methods that have been invented. However, there are several critical obstacles that have limited framework performance, even for state-of-the-art frameworks. Three major challenges include: (i) reliance on large amounts of data, especially in medical imaging, which has been a challenge for effective training; particularly, medical images such as Magnetic Resonance Imaging (MRI) and Computed Tomography (CT) are difficult to collect and refine to create high-quality, consistent datasets; (ii) the inconsistency of samples created by various organ shapes according to slice positions and patient conditions; and (iii) the imbalance between the foreground and background in segmentation tasks, as organ mask areas are small compared to the background. Motivated by that, we propose LoGo - a local-global awareness prototypical segmentation framework to solve these challenges through three synergistic components. The first component is a dual-pathway LoGo encoder responsible for extracting local-grained features and global contextual features under minimal annotation samples; a distillation step is applied to initialize this encoder for better performance in training. Second, we introduce a marginal fairness prototype as a constrained optimizer to address prototype collapses as well as the imbalance problem. The third is a pixel-wise similarity fusion to integrate pixel prototypes with the local patterns, which enables spatially coherent segmentation with sharply preserved anatomical boundaries. Empirical evaluations using abdominal CT and cardiac MRI datasets reveal that LoGo systematically exceeds the performance of current leading models. By attaining a mean Dice score of 90.05\% while maintaining a highly compact parameter set, the model proves both effective and resource-efficient. Furthermore, its ability to generalize across diverse datasets and function with limited labeled data suggests it is a viable candidate for large-scale clinical integration.